\section{Exemples de foncteurs fibres et de points de topos}\etiquette[7]{IV.sec7}\pageoriginale

\subsection{Cas de $\Top(X)$ pour un espace topologique
  $X$}\etiquette[7.1]{IV.subsec7.1}

Soit $x$ un point de $X$. Regardons l'ensemble ponctuel $\{x\}$
comme un espace topologique, et considérons l'inclusion
$$
i_x:\{x\}\longrightarrow X{\quoi}.
$$
En vertu de \hyperref[IV.subsubsec4.1.1]{4.1.1} il définit un morphisme de topos
(défini à isomorphisme unique près)
\begin{equation*}
\Top(i_x):\Top(\{x\})\longrightarrow \Top(X){\quoi}.\tag{7.1.1}
\end{equation*}\etiquette[7.1.1]{IV.eq7.1.1}
D'autre part, on peut identifier $P$ à $\Top(\{x\})$
(\hyperref[IV.subsec2.2]{2.2}), et on trouve ainsi un point $p_x$ de
$\Top(X)$:
\begin{equation*}
p_x:P\longrightarrow \Top(X){\quoi},\tag{7.1.2}
\end{equation*}\etiquette[7.1.2]{IV.eq7.1.2}
défini par $x$ à isomorphisme unique près. Le foncteur fibre associé
est donné par la formule bien connue $[TF]$, résultant d'ailleurs
immédiatement de \hyperrefext[I][I.prop5.1]{5.1}:
\begin{equation*}
p^{\ast}_x(F)=F_x=F_{p_x}=\varinjlim_{U\ni x}F(U){\quoi},\tag{7.1.3}
\end{equation*}\etiquette[7.1.3]{IV.eq7.1.3}
la limite étant prise suivant les voisinages ouverts $U$ de $x$ dans
$X$.

Comme on sait que $\Top(X)$ est isomorphe à $\Top(X_{\sob})$
(\hyperref[IV.subsubsec4.2.1]{4.2.1}), on voit plus généralement que tout point
de $X_{\sob}$, \ie toute partie fermée irréductible $Z$ de $X$,
définit un point de $\Top(X)$, ou encore un foncteur fibre sur
$\Top(X)$, qu'on notera encore $F\mapsto F_Z$. Revenant à sa
définition par la formule \eqhyperref[IV.eq7.1.3]{7.1.3} sur $X_{\sob}$, on trouve
\begin{equation*}
F_Z=\varinjlim_{U\cap Z\neq \phi} F(U){\quoi},\tag{7.1.4}
\end{equation*}\etiquette[7.1.4]{IV.eq7.1.4}
la\pageoriginale limite inductive étant prise suivant les ouverts
$U$ de $X$ qui rencontrent $Z$.

\setcounter{subsubsection}{4}
\subsubsection{}\etiquette[7.1.5]{IV.subsubsec7.1.5}
Il est bien connu que les foncteurs fibres $F\mapsto F_x(x\in X)$
forment déjà une famille conservative. C'est particulièrement
évident sur l'interprétation des faisceaux sur $X$ en termes
d'espaces étalés, la fibre de $F$ en $x$ n'étant alors autre que sa
fibre au sens des espaces fibrés sur $X$. On notera que l'ensemble
d'indices de la famille conservative de foncteurs fibres envisagée
est $X$, donc est $𝒰$-petite.

\subsubsection{}\etiquette[7.1.6]{IV.subsubsec7.1.6}
En vertu de \hyperref[IV.subsubsec4.2.3]{4.2.3}, la catégorie $\Ppoint(\Top(X))$
est équivalente à la catégorie associée à l'ensemble $X_{\sob}$
ordonné par la relation de spécialisation. En d'autres termes: tout
foncteur fibre sur $\Top(X)$ est isomorphe à un foncteur fibre
$F\mapsto F_Z$, où $Z$ est un élément uniquement déterminé de
$X_{\sob}$, \ie une partie fermée irréductible uniquement
déterminée
de $X$; d'autre part, si $Z$, $Z'\in X_{\sob}$, alors l'ensemble
$\Hom(F_Z,F_{Z'})$ est vide ou réduit à un point, ce dernier cas se
présentant si et seulement si $Z$ est une spécialisation de $Z'$
dans $X_{\sob}$, \ie si et seulement si $Z\subset Z'$ comme partie
de $X$. Lorsque $X$ lui-même est sobre, on peut dans ces énoncés
remplacer $X_{\sob}$ par $X$ lui-même.

On notera que le groupe des automorphismes d'un foncteur fibre de
$\Top(X)$ (opposé au groupe des automorphismes du point
correspondant de $\Top(X)$) est toujours réduit au groupe unité
(plus généralement, \hyperref[IV.subsubsec4.2.3]{4.2.3} nous apprend la même
chose pour le groupe des automorphismes de tout morphisme de topos
associés à des espaces topologiques). C'est là un phénomène très
spécial au cas particulier envisagé: voir l'exemple
\hyperref[IV.subsec7.2]{7.2}, ainsi\pageoriginale que VIII 7. Dans le cas des
topos associés à des « problèmes de modules »  (\cf par exemple
[10]), les groupes d'automorphismes des foncteurs fibres ont
d'ailleurs une interprétation remarquable, comme les groupes
d'automorphismes des structures algébriques (sur des corps
algébriquement clos) qu'on se propose de classifier.

\subsubsection{}\etiquette[7.1.7]{IV.subsubsec7.1.7}
On vient de voir comment on peut reconstituer un espace sobre $X$
(ou l'espace sobre $X_{\sob}$ associé à un espace topologique
quelconque), du moins en tant qu'ensemble ordonné par la relation de
spécialisation, à l'aide du topos $\Top(X)$ (qu'il suffit même de
connaître à équivalence près), comme l'ensemble des classes
d'isomorphie de « points »  de $\Top(X)$. D'après ce qui a été
dit dans \hyperref[IV.subsec2.1]{2.1}, on reconstitue également la topologie
de $X$, \ie la famille de ses ouverts, de la façon suivante: pour
tout sous-objet $U$ de l'objet final $e_E$ de $E$, soit $\Pt(U)$
l'ensemble des $x\in X$ tels que $U_x\neq\phi$ (qui s'identifie
d'ailleurs, en vertu de \eqhyperref[IV.eq6.7.2]{6.7.2}, à l'ensemble des classes
d'isomorphie de points du topos induit $E_{/U}$). Alors $U\mapsto
\Pt(U)$ est un isomorphisme d'ensembles ordonnés de l'ensemble des
sous-objets de $e_E$ sur l'ensemble des ouverts de $X$.

\subsubsection{}\etiquette[7.1.8]{IV.subsubsec7.1.8}
La détermination \hyperref[IV.subsubsec7.1.6]{7.1.6} de la catégorie des points
du topos $\Top(X)$ défini par un espace topologique $X$ conduit à
adopter la terminologie suivante pour les points $p$, $p'$ d'un
topos quelconque $E$: on dit que $p$ est une \emph{spécialisation}
de $p'$, ou que $p'$ est une \emph{générisation} de $p$, lorsqu'il
existe un morphisme de $p'$ dans $p$ (au sens de la catégorie
$\Ppoint(E)$ de \hyperref[IV.def6.1]{6.1}). Ces relations sont encore
transitives, mais on trouve facilement grâce à
\hyperref[IV.subsubsec6.8.6]{6.8.6} des exemples où $p$ et $p'$ sont chacun
spécialisation de\pageoriginale l'autre, sans que $p$ et $p'$ soient
isomorphes (ni a fortiori égaux). La catégorie des
générisations
d'un point $p$ du topos $E$, par quoi on entend la catégorie
$\Ppoint(E)_{/p}$, joue à certains égards un rôle analogue à celui
du passage au localisé d'un schéma en un point. Ce rôle sera
précisé
au \sisi{Chap. VI}{\hyperrefext[VI][VI.subsec8.4]{8.4}} avec la construction du topos localisé de $E$ en le
point $p$, dont la catégorie des points est canoniquement
équivalente à la catégorie des générisations de $p$.

\begin{enonce*}[remark]{Exercice 7.1.9}\etiquette[7.1.9]{IV.exer7.1.9}
Soit $E$ un $𝒰$-topos. Prouver que $E$ est équivalent à un
topos de la forme $\Top(X)$, où $X$ est un espace topologique $\in
𝒰$, si et seulement si il satisfait aux deux conditions
suivantes:
\begin{enumerate}
\renewcommand{\theenumi}{\alph{enumi}}
\renewcommand{\labelenumi}{\theenumi)}
\item La famille des sous-objets de l'objet final $e_E$ est
génératrice pour le topos $E$.

\item $E$ a suffisamment de points (\hyperref[IV.prop6.5]{6.5}). (Cette condition n'est pas
superflue: \cf \hyperref[IV.subsec7.4]{7.4}.)
\end{enumerate}

Comparer avec \hyperref[IV.exer7.8]{7.8} c).
\end{enonce*}

\begin{enonce*}[remark]{Exercice 7.1.10}\etiquette[7.1.10]{IV.exer7.1.10}
Soient $X$ un espace topologique muni d'un groupe d'opérateurs $G$
$(X,G\in𝒰)$ et soit $E=\Top(X,G)$ (\hyperref[IV.subsec2.3]{2.3}).
\begin{enumerate}
\renewcommand{\theenumi}{\alph{enumi}}
\renewcommand{\labelenumi}{\theenumi)}
\item Montrer que pour tout objet $(X',G)$ de $E$, le topos induit
$E_{/(X',G)}$ est canoniquement équivalent au topos $\Top(X',G)$
(comparer \hyperref[IV.subsec5.7]{5.7}).

\item Lorsque $G$ opère proprement et librement sur $X'$ (Bourbaki,
Top. Gen. Chap. \hyperrefext[III][III.sec4]{4}), montrer que le morphisme d'espaces à
opérateurs $(X',G)\rightarrow (X'/G,e)$ induit
(\hyperref[IV.subsubsec4.1.2]{4.1.2}) une équivalence de topos
$$
\Top(X',G)\longrightarrow \Top(X'/G){\quoi}.
$$

\item Conclure\pageoriginale de b) qu'il existe un objet $Z$ de $E$
tel que le topos induit $E_{/Z}$ soit équivalent à $\Top(X)$, le
foncteur de localisation $E\rightarrow E_{/Z}$ étant isomorphe au
foncteur « oubli des opérations de $G$ ». (Calquer le
raisonnement de \hyperref[IV.subsubsec5.8.3]{5.8.3})

\item Conclure de c) que tout foncteur fibre sur $E$ est induit, via
le foncteur « oubli des opérations de $G$ », par un foncteur
fibre sur $\Top(X)$, donc définissable par un point de $X_{\sob}$.

\item Déterminer la structure de la catégorie $\Ppoint(E)$ (à
équivalence près) en termes de l'espace à opérateurs $(X_{\sob},G)$.
En conclure en particulier que deux points de $X_{\sob}$ définissent
des foncteurs fibres sur $E$ isomorphes si et seulement si ils sont
conjugués sous l'action de $G$.
\end{enumerate}
\end{enonce*}

\begin{enonce*}[remark]{Exercice 7.1.11}\etiquette[7.1.11]{IV.exer7.1.11}
Soit $X\in 𝒰$ un espace topologique. Prouver que le
$V$-topos $\TOP(X)$ (\hyperref[IV.subsec2.5]{2.5}) a suffisamment de points.
Plus précisément, pour tout objet $X'$ de $\TOP(X)$, et tout $x'\in
X'$, définir un foncteur fibre $F\mapsto F_{X',x'}$ sur $\TOP(X)$,
ne dépendant que du « germe » d'espace $(X',x')$ au-dessus de
$X$, et prouver que cette famille de foncteurs fibres est
conservative (et indexée par un ensemble d'indices
$𝒱$-petit). Montrer, en utilisant un système projectif
filtrant convenable $(X'_i,x'_i)_{i\in I}$, avec $I$ non
$𝒰$-petit, que l'on n'obtient pas de cette façon tous les
foncteurs fibres du $𝒱$-topos $\TOP(X)$. Montrer que
l'ensemble des classes d'isomorphie de tels foncteurs fibres n'est
pas de cardinal $\in 𝒱$.
\end{enonce*}

\subsection{Points d'un topos classifiant
  $B_G$}\etiquette[7.2]{IV.subsec7.2}\pageoriginale

Soient $E$ un $𝒰$-topos, $G$ un Groupe de $E$, d'où un topos
classifiant $B_G$ (\hyperref[IV.subsec2.4]{2.4}). Si $e$ est le Groupe
ponctuel, les morphismes $e\rightarrow G\rightarrow e$ définissent
(\hyperref[IV.subsec4.5]{4.5}) des morphismes de topos (compte tenu que
$B_e\approx E$)
\begin{equation*}
E\longrightarrow B_G \longrightarrow E{\quoi},\tag{7.2.1}
\end{equation*}\etiquette[7.2.1]{IV.eq7.2.1}
dont le composé est isomorphe à $\id_E$, d'où par passage aux
catégories de points des foncteurs
\begin{equation*}
\Ppoint(E)\longrightarrow \Ppoint(B_G)\longrightarrow
\Ppoint(E){\quoi},\tag{7.2.2}
\end{equation*}\etiquette[7.2.2]{IV.eq7.2.2}
dont le composé est isomorphe à l'identité. Utilisons le fait que le
premier morphisme de topos \eqhyperref[IV.eq7.2.1]{7.2.1} s'identifie à un morphisme
de localisation relativement à l'objet $E_G=G$ de $B_G$
(\hyperref[IV.subsec5.8]{5.8}), d'où résulte en vertu de \eqhyperref[IV.eq6.7.2]{6.7.2} que le
premier foncteur \eqhyperref[IV.eq7.2.2]{7.2.2} s'identifie au foncteur naturel « d'inclusion »
\begin{equation*}
\Ppoint(B_G)_{/\widehat{E}_G}\longrightarrow
\Ppoint(B_G){\quoi},\tag{7.2.3}
\end{equation*}\etiquette[7.2.3]{IV.eq7.2.3}
où $\widehat{E}_G$ désigne le préfaisceau $p'\mapsto (E_G)_{p'}$ sur
$\Ppoint(B_G)$. Comme $E_G$ couvre évidemment l'objet final de
$B_G$, ses fibres $(E_G)_p$ sont non vides, d'où résulte que
\eqhyperref[IV.eq7.2.3]{7.2.3} est essentiellement surjectif, ce qui signifie que
\emph{tout foncteur fibre sur $B_G$ est induit (à isomorphisme près)
par un foncteur fibre de $E$, via le foncteur « oubli des
opérations de $G$ » $B_G\rightarrow E$.} On en conclut en
particulier que \emph{si $E$ a suffisamment de points (\resp une
petite famille conservative de points)} il en est de même de $B_G$.

\begin{enonce*}[remark]{Remarque 7.2.4}\etiquette[7.2.4]{IV.rem7.2.4}
On peut aller plus loin et déterminer (à équivalence près) la
structure de la catégorie $C'=\Ppoint(B_G)$ en termes de la
catégorie $C=\Ppoint(E)$ et du préfaisceau en groupes
$$
\widehat{G}:p\longmapsto G_p: C^{\circ}\longrightarrow
(\text{Groupes})
$$
sur\pageoriginale celle-ci. Définissons en effet une nouvelle
catégorie $C_1$, ayant mêmes objets que $C$, et telle que pour deux
objets $p$, $q$ de $C$, on ait
$$
\Hom_{C_1}(p,q)=\Hom_C(p,q)\times \widehat{G}(p){\quoi},
$$
la composition des flèches dans $C_1$ étant induite par celle de
$C$, et par la loi de groupe du préfaisceau $\widehat{G}$. La donnée
d'un foncteur $\varphi_1:C_1\rightarrow D$ dans une catégorie
quelconque $D$ équivaut alors à la donnée d'un foncteur
$\varphi:C\rightarrow D$, muni d'une opération du préfaisceau
$\widehat{G}$ sur ce dernier (\ie la donnée, pour tout $p\in \ob C$,
d'une opération de $\widehat{G}(p)$ sur $\varphi(p)$, satisfaisant à
une condition de fonctorialité évidente pour $p$ variable).
Utilisant cette observation, on définit un foncteur canonique
\begin{equation*}
\varphi_1:C_1\longrightarrow C'{\quoi},\tag{7.2.4.1}
\end{equation*}\etiquette[7.2.4.1]{IV.eq7.2.4.1}
correspondant au foncteur $\varphi:C\rightarrow C'$ de
\eqhyperref[IV.eq7.2.2]{7.2.2}, associant à tout point $p$ de $E$ le point
$\varphi(p)$ induit sur $B_G$, avec les opérations naturelles de
$\widehat{G}(p)=G_p$ sur ce dernier, provenant des opérations de
$G_p$ sur les $X_p$ lorsque $X$ parcourt $B_G$. On laisse au lecteur
le soin de prouver que \eqhyperref[IV.eq7.2.4.1]{7.2.4.1} est une équivalence de
catégories, en utilisant la structure \eqhyperref[IV.eq7.2.3]{7.2.3} du premier
foncteur de \eqhyperref[IV.eq7.2.2]{7.2.2}, et en notant que le foncteur naturel
d'inclusion $C\rightarrow C_1$ admet une structure analogue,
relativement à l'image inverse $P_1$ du préfaisceau
$P'=\widehat{E}_G$ sur $C'$.
\end{enonce*}

\setcounter{subsubsection}{4}
\subsubsection{}\etiquette[7.2.5]{IV.subsubsec7.2.5}
On\pageoriginale conclut en particulier que les \emph{classes
d'isomorphie de points de $B_G$ correspondent exactement aux classes
d'isomorphie de points de $E$: tout point du topos classifiant $E_G$
est induit, à isomorphisme} non unique \emph{près, par un point de
$E$}. En particulier, lorsque $E$ est le topos ponctuel $P$, donc
que $G$ est un groupe ordinaire, alors la catégorie $\Ppoint(B_G)$
est un groupoïde connexe à groupe fondamental $G$: tout foncteur
fibre sur $B_G$ est isomorphe (de façon non canonique) au
foncteur $\omega$ « oubli des opérations de $G$ », et le monoïde
des endomorphismes de ce dernier foncteur est le groupe $G$ (donc
tout endomorphisme de $\omega$ est un automorphisme).

\begin{enonce*}[remark]{Exercice 7.2.6}\etiquette[7.2.6]{IV.exer7.2.6}
\begin{enumerate}
\renewcommand{\theenumi}{\alph{enumi}}
\renewcommand{\labelenumi}{\theenumi)}
\item Soit (Tors) la catégorie des couples $(\Gamma,P)\in 𝒰$,
avec $\Gamma$ un groupe et $P$ un torseur à droite sous $\Gamma$. La
foncteur $(\Gamma,P)\mapsto \Gamma$
\begin{equation*}
(\Tors)\longrightarrow (\Groupes)\tag{7.2.6.1}
\end{equation*}\etiquette[7.2.6.1]{IV.eq7.2.6.1}
est un foncteur cofibrant. Avec les notations de \hyperref[IV.rem7.2.4]{7.2.4},
considérons le foncteur
$$
\widehat{G}:C^{\circ}\longrightarrow (\Groupes){\quoi},
$$
et soit $D'$ la catégorie cofibrée sur $C^{\circ}$ image inverse de
la catégorie cofibrée \eqhyperref[IV.eq7.2.6.1]{7.2.6.1}, $D$ la
catégorie fibrés
correspondants sur $C$, ayant même catégories fibres que $D'$, de
sorte que pour $p\in \ob C$, on ait
$$
D_p=\mbox{ catégorie des $G_p$-torseurs à droite}{\quoi}.
$$
Construire des foncteurs canoniques
\begin{equation*}
D\longrightarrow C'=\Ppoint(B_G)\quav C'\longrightarrow
D{\quoi},\tag{7.2.6.2}
\end{equation*}\etiquette[7.2.6.2]{IV.eq7.2.6.2}
\emph{quasi-inverses l'un de l'autre.} En particulier, en conclure
que la catégorie des points de $B_G$ au-dessus d'un point donné $p$
de $E$ est canoniquement équivalente à la catégorie des torseurs à
droite sous le groupe $G_p$.

\item Soit\pageoriginale $𝒢=(G_i)_{i\in I}$ un système
projectif strict de Groupes du $𝒰$-topos $E$ ($I\in
𝒰$ un ensemble préordonné filtrant), d'où un topos
classifiant $B_{𝒰}$, en calquant la définition de
\hyperref[IV.subsubsec2.7.1]{2.7.1}. Déterminer la catégorie des points de
$B_{𝒢}$, en calquant a). (NB. On remplacera les catégories
$(\Tors)$ et $(\Groupes)$ par les catégories de systèmes projectifs
indexés par $I$ de ces catégories.) En conclure que si $I$ admet un
ensemble cofinal dénombrable, alors tout foncteur fibre sur
$B_{𝒢}$ est isomorphe à un foncteur induit par un foncteur
fibre de $E$ (via le foncteur « oubli des opérations de $G$ »),
ce dernier étant déterminé à isomorphisme non unique près.

\item Prenant pour $E$ le topos ponctuel, de sorte que $𝒢$ est
un pro-groupe ordinaire, montrer que la conclusion de b) est
également valable si $I$ est quelconque, mais en revanche
$𝒢$ est profini (\ie les $G_i$ sont finis): tout foncteur
fibre est isomorphe au foncteur « oubli des opérations de $G$ ».

\item Prenant toujours le cas où $E$ est le topos ponctuel, donner un
exemple d'un foncteur fibre sur $B_{𝒢}$, pour un système
projectif convenable $𝒢=(G_i)_{i\in I}$, qui n'est pas
isomorphe au foncteur « oubli des opérations de $G$ ».

\item Considérons $𝒢=(G_i)$ comme un Groupe du topos
$\widehat{I}$, et soit $𝒯$ une \emph{gerbe} sur $I$ de
\emph{lien} $G$ [\hyperref[IV.sec3]{3}]. Montrer comment on
peut « tordre » le topos classifiant $B_G$ à l'aide de la gerbe $T$, pour
obtenir un topos $B^T_{𝒢}$, limite inductive de
sous-catégories $B^T(i)$ équivalentes (non canoniquement) aux topos
classifiants $B_{G_i}$. Montrer que le topos $B^T_{𝒢}$ admet
un foncteur fibre si et seulement si la gerbe $T$ est
« neutre », en établissant une\pageoriginale équivalence de catégories
entre la catégorie des foncteurs fibres de $B^T_{𝒢}$ et la
catégorie des sections de $T$ sur $I$. En conclure, si les $G_i$
sont commutatifs, de sorte que la classification des gerbes sur $I$
de lien $𝒢$ se fait par
$\displaystyle{\varprojlim_I}^{(2)}G_i=H^{2}(I,𝒢)$
(\loccit), un exemple d'un topos (non « vide »
(\hyperref[IV.subsec2.2]{2.2})) de la forme $B^T_{𝒢}$ qui n'a pas de
points. (Prendre un exemple où
$\displaystyle{\varprojlim_I}^{(2)}G_i\neq 0$.)
\end{enumerate}
\end{enonce*}

\subsection{Points des topos $\widehat{C}$; exemples de
$𝒰$-topos $\widehat{C}$ dont la catégorie
des points ne soit pas équivalente à une petite
catégorie}\etiquette[7.3]{IV.subsec7.3}

Soient $E$ un $𝒰$-topos, $(\varphi_i)_{i\in I}$ une famille
filtrante de foncteurs fibres sur $E$. Il résulte des propriétés
d'exactitude des foncteurs $\varinjlim$ (\sisi{I\quad et}{\hyperrefext[I][I.prop2.8]{2.8}}) que le foncteur
$\varphi=\varinjlim\varphi_i$ est également un foncteur fibre:
\emph{toute limite inductive filtrante de foncteurs fibres est un
foncteur fibre.} Par suite, la catégorie $\Ffib(E)$ admet des
$𝒰$-limites inductives filtrantes (et le foncteur
d'inclusion $\Ffib(E)\rightarrow \Hhom(E,(𝒰\text{-}\Ens))$ y
commute); en d'autres termes, la catégorie $\Ppoint(E)$ admet des
$𝒰$-limites projectives filtrantes. Ce fait a déjà été
utilisé dans exercice \sisi{\hyperref[IV.exer7.1.10]{7.1.10}}{\hyperref[IV.exer7.1.11]{7.1.11}} pour donner un exemple de
« grosses »  catégories de foncteurs fibres.

On peut construire des exemples nettement plus simples, avec des
topos de la forme $E=\widehat{C}$, $C\in𝒰$, en utilisant
\eqhyperref[IV.eq6.8.6.1]{6.8.6.1}. Ceci nous donne aussitôt des exemples où
$\Ffib(\widehat{C})$ n'est pas équivalente à une catégorie
$\in𝒰$ \ie où le cardinal de l'ensemble des classes
d'isomorphie d'objets de cette $𝒰$-catégorie n'est pas
$\in𝒰$. Ceci signifie que la catégorie $\Pro-(C)$ n'est pas
équivalente à une catégorie $\in 𝒰$. Il suffit par exemple
de prendre pour $D=C^{\circ}$ la catégorie des ensembles finis de la
forme $[0,n]$, où $n\geq 0$ est un\pageoriginale entier, auquel cas
$\Ind(D)\simeq \Pro(C)^{\circ}$ est équivalente à la catégorie des
ensembles non vides. Le topos $E$ est dans ce cas la catégorie bien
connue des \emph{ensembles cosimpliciaux}. On pourrait aussi prendre
pour $C$ la catégorie des ensembles finis non vides, on trouve que
la catégorie des points sur le topos $\widehat{C}$ des ensembles
simpliciaux est équivalente à la catégorie des ensembles profinis,
ou encore à la catégorie des espaces compacts totalement
discontinus.

\subsection{Topos non vides sans points}\etiquette[7.4]{IV.subsec7.4}

L'exemple suivant est dû à $P$. Deligne. (Pour un autre exemple, \cf
\hyperref[IV.exer7.2.6]{7.2.6} e).) On prend un espace compact $K$ muni d'une
mesure $\mu$, et l'ensemble ordonné $U$ des parties mesurables de
$K$ à ensemble de mesure nulle près. On fait de $U$ une catégorie
$\underline{U}$ telle que $\ob \underline{U}=U$, les morphismes de
$\underline{U}$ étant les « morphismes d'inclusion »  entre éléments
de $U$. On fait de $\underline{U}$ un site en prenant la prétopologie
pour laquelle $\Cov(E)$ (pour $E\in U$) est formé des familles
\emph{dénombrables} d'éléments $E_i$ de $E$ majorés par $E$, telles
que $E$ soit la réunion des $E_i$ à ensemble de mesure nulle près.
On en déduit un Topos $\Top(\mu)=\underline{U}^{\sim}$, admettant
l'ensemble des sous-objets de l'objet final comme famille
génératrice (lequel topos semble avoir échappé à l'attention des
probabilistes). Ce topos est un « topos vide »
(\hyperref[IV.subsec2.2]{2.2}) si et seulement si $\mu=0$. D'autre part, la
catégorie des points de ce topos est équivalent à la catégorie
discrète définie par l'ensemble des points $x\in K$ tels que
$\mu(\{x\})\neq 0$ (démonstration au lecteur). Elle est donc vide si
$K$ n'admet pas de tels points, par exemple si $K$ est le segment
unité de la droite, avec la mesure induite par la mesure de
Lebesgue.

\begin{enonce*}[remark]{Exercice 7.5}\etiquette[7.5]{IV.exer7.5}
(\emph{catégories\pageoriginale \sisi{Karoubiennes}{karoubiennes} et morphismes de
topos})\footnote[1]{Comparer \hyperrefext[I][I.exer8.7.8]{8.7.8}.}
\begin{enumerate}
\renewcommand{\theenumi}{\alph{enumi}}
\renewcommand{\labelenumi}{\theenumi)}
\item Soient $C$ une catégorie, $X$ un objet de $C$, $p$ un
endomorphisme de $X$. On dit que $p$ est un \emph{projecteur} si
$p^2=p$. Prouver que si $p$ est un projecteur, pour que
$\Ker(\id_X,p)$ soit représentable, il faut et il suffit que
$\Coker(\id_X,p)$ le soit, et que les deux objets de $C$ ainsi
obtenus sont canoniquement isomorphes. On dira alors que la
projecteur $p$ \emph{admet une image,} et $\Ker(\id_X,p)\simeq
\Coker(\id_X,p)$ est appelé l'\emph{image du projecteur $p$}; on
l'identifie suivant le contexte à un sous-objet ou à un objet
quotient de $C$. Un objet isomorphe à un $\text{Im\,} p$, pour un
projecteur convenable $p$ dans $X$, est appelé un \emph{facteur
direct de l'objet $X$ de $C$.} On dit que $C$ est une
\emph{catégorie avec facteurs directs}, ou une \emph{catégorie
\sisi{Karoubienne}{karoubienne},} si tout projecteur dans un objet de $C$ admet une
image. Si $F:C\rightarrow C'$ est un foncteur qui commute aux noyaux
ou aux conoyaux, alors $F$ transforme un projecteur admettant une
image en un projecteur admettant une image\nde{Pas besoin de conditions sur
$F$ ici, tous les foncteurs préservent les projecteurs admettant des images.}.

\item Montrer que pour toute catégorie $C$, on peut trouver un
foncteur $\varphi:C\rightarrow \kar (C)$ de $C$ dans une catégorie
karoubienne (déterminée à équivalence près), tel que pour toute
catégorie karoubienne $C'$, le foncteur
$$
f\mapsto f\circ \varphi:\Hhom(\kar(C),C')\longrightarrow \Hhom(C,C')
$$
soit une équivalence de catégories; on appellera $\kar(C)$
l'\emph{enveloppe de Karoubi} de la catégorie $C$. (Hint: prendre
pour $\ob \kar (C)$ l'ensemble des couples $(X,p)$, avec $X\in \ob
C$ et $p$ un projecteur dans $X$, et pour $\Hom((X,p),(Y,q))$ la
partie de $\Hom(X,Y)$ formée des $f:X\rightarrow Y$ tels que
$f=qfp.$) Montrer que $\varphi$ est pleinement fidèle.

\item Soit $F:E\rightarrow E'$ un foncteur d'une catégorie dans une
autre. Montrer que si $F$ commute à un certain type de limites
inductives ou projectives,\pageoriginale il en est de même de tout
facteur direct de $F$. En particulier, si $E$ et $E'$ sont des
$𝒰$-topos et si $F$ est un foncteur image inverse pour un
morphisme de topos $E'\rightarrow E$, alors il en est de même de
tout facteur direct de $F$; par suite, la catégorie $\Hhomtop(E',E)$
est karoubienne, et en particulier la catégorie $\Ppoint(E)$ est
karoubienne.

\item Montrer que pour toute $𝒰$-catégorie $C$, la catégorie
$\Ind(C)$ est karoubienne (où $\Ind(C)$ est formée avec les systèmes
inductifs de $C$ indexés par un ensemble préordonné filtrant $I\in
𝒰$). (Si $C\in 𝒰$, utiliser par exemple le fait
(\hyperref[IV.subsec7.3]{7.3}) que $\Ind(C)$ est équivalente à la catégorie
$\Ffib((C^{\circ})^{\widehat{\phantom{,}}})=\Ppoint((C^{\circ})^{\widehat{\phantom{,}}})^{\circ}.$)
En conclure une autre construction de $\kar(C)$ comme sous-catégorie
pleine de $\Ind(C)$ formée des images dans $\Ind(C)$ des projecteurs
d'objets $X$ de $C$.
\end{enumerate}
\end{enonce*}

\begin{enonce*}[remark]{Exercice 7.6}\etiquette[7.6]{IV.exer7.6}
(\emph{Morphismes essentiels de topos, points essentiels}).
\begin{enumerate}
\renewcommand{\theenumi}{\alph{enumi}}
\renewcommand{\labelenumi}{\theenumi)}
\item Soit $f:E\rightarrow F$ un morphisme de topos. Montrer que les
conditions suivantes sont équivalentes:
\begin{itemize}
\item[i)] $f_{!}$ existe, \ie $f^{\ast}$ admet un adjoint à gauche.

\item[ii)] $f^{\ast}$ commute aux $𝒰$-limites projectives.

\item[${\rm ii')}$] $f^{\ast}$ commute aux $𝒰$-produits.
\end{itemize}
On dira alors que $f$ est un \emph{morphisme essentiel} du topos $E$
dans le topos $E'$.

Montrer que si $f$ satisfait à ces conditions, il en est de même de
tout facteur direct de $f$. (Utiliser \hyperref[IV.rem1.8]{1.8} et
\hyperref[IV.exer7.5]{7.5} c).)

\item Soit $E$ un topos. On appelle \emph{point essentiel du topos}
$E$ tout point $p:P\rightarrow E$ de $E$ tel que $p_{!}$ existe.
Montrer que si $E=\Top(X)$, $X$ espace topologique, et si $p$ est le
point de $E$ défini par un $x\in X$, alors $p$ est
essentiel\pageoriginale si et seulement si $x$ admet un plus petit
voisinage ouvert $U$ (ce qui signifie, si les points de $X$ sont
fermés, que $x$ est un \emph{point isolé} de $X$). Pour qu'un
morphisme de topos $f:E\rightarrow F$ soit essentiel (a)), il faut
que $f$ transforme points essentiels en points essentiels, et cette
condition est également suffisante lorsque $E$ admet suffisamment de
points essentiels (\ie si la famille de ces points est conservative)
(\cf d)).

\item Pour qu'un point $p$ du topos $E$ soit essentiel, il faut et il
suffit que le foncteur fibre associé soit représentable par un objet
$X$ de $E$. Pour que le foncteur covariant $\varphi:E\rightarrow
(\Ens)$ représenté par un objet donné $X$ de $E$ soit un foncteur
fibre (\ie définisse un point, nécessairement essentiel, de $E$), il
faut et il suffit que $X$ soit connexe non vide
(\hyperref[IV.subsubsec4.3.5]{4.3.5}), et projectif (\ie que le foncteur
$\varphi$ transforme épimorphismes en épimorphismes). (Hint: montrer
d'abord que pour que $\varphi$ commute aux sommes, il faut et il
suffit que $X$ soit connexe non vide, puis utiliser le critère
\hyperref[IV.prop4.9.4]{4.9.4}.) En conclure une équivalence entre la catégorie
$\Ppointess(E)$ et la sous-catégorie pleine $P$ de $E$ formée des
objets connexes non vides projectifs.

\item Montrer que la topologie de $P$ induite par celle de $E$ est la
topologie chaotique. En conclure l'équivalence des conditions
suivantes sur $E$ (due à J.E ROOS [12 c), prop. 1]): (i) La famille
des points essentiels de $E$ est conservative; (ii) La
sous-catégorie pleine $P$ de $E$ formée des objets connexes-non
vides projectifs est génératrice; (iii)~$E$ est équivalent à un
topos de la forme $\widehat{C}$, où $C$ est une catégorie
équivalente à une catégorie $\in 𝒰$ (ou, si on préfère,
$C\in 𝒰$). (Utiliser $c$) et le fait que si $P$ est
génératrice, le foncteur naturel $E\rightarrow P^{\sim}$ est une
équivalence de\pageoriginale catégories.)

\item La catégorie des points essentiels d'un topos $E$ est
karoubienne (\cf (\hyperref[IV.exer7.5]{7.5} c)). Soit $C$ une catégorie
équivalente à une catégorie $\in 𝒰$. Prouver que le foncteur
pleinement fidèle canonique (\hyperref[IV.subsubsec4.6.2]{4.6.2})
$$
C\longrightarrow \Ppoint(\widehat{C})
$$
se factorise par un foncteur
$$
C\longrightarrow \Ppointess(\widehat{C}){\quoi},
$$
qui fait de $\Ppointess(\widehat{C})$ une \emph{enveloppe de
Karoubi} (\hyperref[IV.exer7.5]{7.5} b)) \emph{de} $C$. En particulier, ce
foncteur est une équivalence de catégories si et seulement si la
catégorie $C$ est karoubienne. (Hint: utilisant $c$), prouver que
tout point isolé de $\widehat{C}$ est isomorphe à un facteur direct
d'un $X\in \ob C$.)

\item Soient $C$ et $C'$ deux catégories équivalentes à des catégories
$\in 𝒰$. Prouver que le foncteur pleinement fidèle canonique
(\hyperref[IV.subsubsec4.6.2]{4.6.2})
$$
\Hhom(C,C')\longrightarrow \Hhomtop(\widehat{C},\widehat{C}')
$$
prend ses valeurs dans la sous-catégorie pleine
$\Hhomtopess(\widehat{C},\widehat{C}')$ des morphismes essentiels
\emph{de topos} (a)), et que le foncteur induit
$$
\Hhom(C,C')\longrightarrow \Hhomtopess(\widehat{C},\widehat{C}')
$$
est une équivalence de catégories si et seulement si $C$ est vide ou
$C'$ est karoubienne. (Utiliser $g$) ci-dessous.)

\item Soient $C$ une catégorie équivalente à une catégorie $\in
𝒰$, et $E$ un topos. Définir une équivalence entre la
catégorie $\Hhomtopess(\widehat{C},E)$ des morphismes essentiels de
topos $\widehat{C}\rightarrow E$, et la catégorie $\Hhom(C,
\Ppointess(E))$. (Utiliser $b$)\pageoriginale et le fait que
$\widehat{C}$ a suffisamment de points essentiels).

\item Soit $C$ une catégorie finie. Prouver que tout point de
$\widehat{C}$ est essentiel, et que $\Ppoint(\widehat{C})$ est
équivalente à une catégorie finie. (Noter que l'enveloppe de Karoubi
de $C$ est finie, et utilisant \eqhyperref[IV.eq6.8.6.1]{6.8.6.1}, se ramener à prouver
que si $C$ est une catégorie karoubienne finie, alors le foncteur
canonique $C\rightarrow \Pro(C)$ est une équivalence de catégories).
Utilisant b), et conclure que \emph{tout} morphisme de topos
$\widehat{C}'\rightarrow \widehat{C}$ est essentiel, et
$\Hhom(C,C')\rightarrow \Hom(\widehat{C},\widehat{C}')$ est une
équivalence si et seulement si $C$ est vide ou $C'$ karoubienne.

\item Soit $X$ un espace topologique, $f:\Top(X)\rightarrow P$ le
morphisme de topos déduit de l'application continue de $X$ dans
l'espace ponctuel. Montrer que $f$ est essentiel si et seulement si
l'espace $X$ est \emph{localement connexe}, \ie satisfait la
condition suivante: pour tout $x\in X$ et tout voisinage ouvert $U$
de $x$, la composante connexe de $x$ dans $U$ est un voisinage de
$x$, ou encore: pour tout ouvert $U$ de $X$, les composantes
connexes de $U$ sont ouvertes. (\cf \hyperref[IV.exer8.7]{8.7} l) pour
généralisation).
\end{enumerate}
\end{enonce*}

\begin{enonce*}[remark]{Exercice 7.7}\etiquette[7.7]{IV.exer7.7}
(\emph{Points inhabituels d'un topos classifiant})\footnote[1]{\cf
aussi \hyperref[IV.exer7.2.6]{7.2.6} d).}.
\begin{enumerate}
\renewcommand{\theenumi}{\alph{enumi}}
\renewcommand{\labelenumi}{\theenumi)}
\item Soit $G$ un monoïde. Pour que $G$ contienne un élément $g$ tel
que $g\neq e$, $g^2=g$, il faut et il suffit que le topos
classifiant $B_G$ admette un point essentiel qui n'est pas isomorphe
au point banal (correspondant au foncteur fibre « oubli des
opérations de $G$ »). (Utiliser \hyperref[IV.exer7.6]{7.6} e)).

\item Soit $G$ le monoïde additif des entiers $\geq 0$. Montrer que
tout point essentiel du topos classifiant $B_G$ est isomorphe au
point banal. Construire un point de $B_G$ qui n'est pas isomorphe au
point banal. Montrer que la catégorie $\Ppoint(B_G)$ n'est pas
équivalente à une catégorie $\in 𝒰$.
\end{enumerate}
\end{enonce*}

\begin{enonce*}[remark]{Exercice 7.8}\etiquette[7.8]{IV.exer7.8}
(\emph{Topologie sur $\Point(E)$, et topos associés aux ensembles
ordonnés}).
\begin{enumerate}
\renewcommand{\theenumi}{\alph{enumi}}
\renewcommand{\labelenumi}{\theenumi)}
\item Soit $E$ un topos. Désignons par $\Point(E)$ l'ensemble des
classes, à isomorphisme près, de points de $E$. Pour tout ouvert $U$
de $E$, soit $U'$ l'ensemble\pageoriginale des classes de points $p$
de $E$ tels que $U_p\neq \phi$. Montrer que $U\mapsto U'$ est une
application croissante de l'ensemble ordonné $\Ouv(E)$ des ouverts
de $E$ dans l'ensemble des parties de $X=\Point(E)$, et que cette
application commute aux Inf finis aux $\Sup$ quelconques. En
conclure que l'ensemble des parties de $X$ de la forme $U'$ définit
sur $X$ une topologie (dite \emph{topologie canonique sur l'ensemble
des classes de points du topos $E$}). Montrer que si $E$ a
suffisamment de points, l'application $U\mapsto U'$ est un
isomorphisme de l'ensemble ordonné $\Ouv(E)$ avec l'ensemble ordonné
$\Ouv(X)$.

\item Soit $\Oo\in 𝒰$ un ensemble ordonné admettant des Inf
fins et des $\Sup$ quelconques, qui définit donc une catégorie
$\cat(\Oo)$ ayant des produits finis et des sommes quelconques. Nous
dirons qu'une famille de morphismes $U_i\rightarrow U$ de même but
$U$ est couvrante si $U$ est le $\Sup$ des $U_i$. Supposant que dans
$\cat(\Oo)$ les sommes sont universelles \ie que dans $\Oo$ les
$\Sup$ quelconques commutent aux Inf, on définit ainsi sur
$\cat(\Oo)$ une topologie qui en fait un site, d'où un topos
$\cat(\Oo)^{\sim}$, noté aussi simplement $\Oo^{\sim}$. Montrer que
$\Oo$ est canoniquement isomorphe à $\Ouv(\Oo^{\sim})$.

Montrer que la catégorie $\Hhomtop(E,\Oo^{\sim})$ est équivalente à
la catégorie associée à l'ensemble ordonné des morphismes
d'ensembles ordonnés $\Oo\rightarrow \Ouv(E)$ qui commutent aux Inf
finis et aux $\Sup$ quelconques. (Utiliser le critère
\hyperref[IV.prop4.9.4]{4.9.4}.) Si $\Oo'$ est un ensemble ordonné satisfaisant
aux mêmes conditions que $\Oo$, conclure que $\Hhomtop(\Oo^{\sim},
{\Oo'}^{\sim})$ est équivalente à la catégorie associée à l'ensemble
ordonné de tous les morphismes d'ensembles ordonnés
$\Oo\rightarrow{\Oo}'$ qui commutent aux Inf finis et aux $\Sup$
quelconques.

\item Montrer\pageoriginale que pour que le topos $E$ soit équivalent
à un topos de la forme $\Oo^{\sim}$, avec $\Oo$ comme dans b), il faut
et il suffit que la famille des ouverts de $E$ soit génératrice dans
$E$. Pour qu'il soit équivalent à un $\Top(X)$ pour $X\in
𝒰$, il faut et il suffit qu'il satisfasse à la condition
précédente et qu'il ait suffisamment de points. Supposant réalisée
la première condition, exprimer la seconde en termes de la structure
de l'ensemble ordonné $\Oo=\Ouv(E)$, en utilisant la dernière
assertion de b).

\item Définir un morphisme de topos canonique
$$
E\longrightarrow \Ouv(E)^{\sim}{\quoi},
$$
qui soit universel (à équivalence près) pour les morphismes de $E$
dans des topos engendrés par leurs ouverts (\cf c)).

\item Soit $X'$ un \emph{petit} sous-ensemble de $X=\Point(E)$, qu'on
munit de la topologie induit par celle de $X$. Définir un morphisme
canonique de topos
$$
\Ouv(E)^{\sim}\longrightarrow \Top(X'){\quoi},
$$
qui est une équivalence lorsque la famille $X'$ de points de $E$ est
conservative. En conclure alors un morphisme canonique de topos
$$
E\longrightarrow \Top(X'){\quoi}.
$$
Donner une caractérisation universelle (à équivalence près) de ce
morphisme de topos pour les morphismes du topos $E$ dans des topos
de la forme $\Top(Y)$. (Noter à ce propos qu'à équivalence près,
$\Top(X')$ ne dépend pas de la petite famille conservative de
foncteurs fibres choisie, ou ce qui revient au même
(\hyperref[IV.subsec4.2]{4.2}), que $X'_{\sob}$ ne dépend pas à homéomorphisme
près du choix d'une telle famille).

\item Supposons\pageoriginale que la catégorie $\Ppoint(E)$ ne
soit pas équivalente à une petite catégorie (\cf
\hyperref[IV.subsec7.3]{7.3}), et que $E$ admette une petite famille
conservative de foncteurs fibres. Montrer que si une telle famille
est prise assez grande, la partie $X'$ de $\Ppoint(E)$ qu'elle
définit n'est pas un espace topologique sobre pour la topologie
induite.

\item Pour les relations avec l'exercice \hyperref[IV.exer7.6]{7.6}, \cf \hyperref[IV.sec8]{8}.
\end{enumerate}
\end{enonce*}

\section{Localisation. Ouverts d'un topos}\etiquette[8]{IV.sec8}

\subsection{}\etiquette[8.1]{IV.subsec8.1}
Soit $C$ une catégorie, et $T(X)$ une relation faisant intervenir un
objet $X$ de $C$. On dit que la relation $T(X)$ est \emph{stable par
changement de base} (en $X$), si pour tout morphisme $X'\rightarrow
X$ de $C$, la relation $T(X)$ implique $T(X')$. Il revient au même
de dire que la sous-catégorie pleine $R$ de $C$, formée des objets
$X$ de $C$ tels qu'on ait $T(X)$, est un \emph{crible}
(\hyperrefext[I][I.def4.1]{4.1}).

\begin{enonce*}[remark]{Remarque 8.1.1}\etiquette[8.1.1]{IV.rem8.1.1}
Supposons que $C$ soit une $𝒰$-catégorie (\hyperrefext[I][I.def1.1]{1.1}),
de sorte que $C$ peut être identifié à une sous-catégorie pleine de
$\widehat{C}$ (\hyperrefext[I][I.prop1.4]{1.4}). Il est parfois commode d'étendre la
définition de la relation $T$ dans $\ob C$ en une relation
$\widehat{T}$ portant sur un objet $F$ de $\widehat{C}$, en
désignant par $\widehat{T}(F)$ la relation: pour toute flèche
$X\rightarrow F$ dans $\widehat{C}$, avec $X\in \ob C$, on a $T(X)$.
La relation $\widehat{T}(F)$ est évidemment encore stable par
changement de base en $F$. Désignant encore par $R$ le sous-objet de
l'objet final $e$ de $\widehat{C}$ défini par le crible $R$ de $C$
(\hyperrefext[I][I.rem4.2.1]{4.2.1}), la relation $\widehat{T}(F)$ peut aussi
s'exprimer par $\Hom_C^{\widehat{\phantom{,}}}(F,R)\neq ∅$, \ie
elle signifie que l'unique morphisme $F\rightarrow e$ se factorise
par le sous-objet $R$ de $e$.
\end{enonce*}

\subsection{}\etiquette[8.2]{IV.subsec8.2}
Supposons maintenant que $C$ soit un \emph{site}. Une relation
$T(X)$ en un argument $X\in \ob C$ est dite \emph{stable par
descente} si pour toute famille couvrante\pageoriginale
$(X_i:X_i\rightarrow X)_{i\in I}$ d'un objet $X$ de $C$, la relation
« $T(X_i)$ pour tout $i\in I$ »  implique la relation $T(X)$. On
dit que $T$ est une \emph{relation de nature locale} si elle est
stable par changement de base (\hyperref[IV.subsec8.1]{8.1}) et stable par
descente. Lorsque $C$ est une $𝒰$-catégorie, et que $T$ est
déjà supposée stable par changement de base, \ie qu'elle est
définissable par un crible $R$ de $C$, qu'on identifiera à un
sous-préfaisceau du préfaisceau final sur $C$, on voit aussitôt que
$C$ est de nature locale si et seulement si $R$ est un
\emph{faisceau}. Ainsi, \emph{les relations de nature locale en un
argument $X\in \ob C$ correspondent exactement aux sous-faisceaux du
faisceau final de $C$, \ie aux sous-objets de l'objet final du topos
$C^{\sim}$.} Elles ne dépendent donc essentiellement que du topos
$C^{\sim}$ défini par le site $C$ (comme toutes les notions
importantes associées à un site !). En termes du sous-faisceau $R$
du faisceau final, la relation $T(X)$ s'exprime comme
$\Hom(\epsilon(X),R)\neq ∅$. Elle se prolonge canoniquement en la
relation $T^{\sim}(F):\Hom(F,R)\neq ∅$ sur le topos $C^{\sim}$.

\begin{enonce*}[remark]{Remarque 8.2.1}\etiquette[8.2.1]{IV.rem8.2.1}
Avec la notation introduite dans \hyperref[IV.rem8.1.1]{8.1.1}, pour un
préfaisceau $F$ sur $C$, comme $\Hom(F,R)\simeq
\Hom(\underline{a}(F),R)$, la relation $\widehat{T}(F)$ équivaut à la
relation $\widehat{T}(\underline{a}(F))$.
\end{enonce*}

\begin{enonce*}{Définition 8.3}\etiquette[8.3]{IV.def8.3}
On appelle ouvert d'un $𝒰$-topos $E$ tout sous-objet de
l'objet final $e_E$ de $E$; si $X$ est un objet de $E$, on appelle
parfois ouvert de $X$ tout ouvert du topos induit $E_{/X}$, \ie tout
sous-objet de $X$.

On appelle ouvert d'un $𝒰$-site $C$ un ouvert du
$𝒰$-topos associé $C^{\sim}$.
\end{enonce*}

On notera que, les objets finaux de $E$ étant canoniquement
isomorphes, l'ambiguïté introduite dans \hyperref[IV.def8.3]{8.3} par le choix
de $e_E$ inoffensive; on\pageoriginale peut aussi, de façon plus
intrinsèque mais moins maniable du point de vue pratique, définir
les ouverts de $E$ comme étant les sous-catégories pleines $R$ de
$E$ qui sont telles que la relation $F\in \ob R$ pour un objet $F$
de $E$ soit une relation de nature locale. De même, les ouverts d'un
$𝒰$-site $C$ correspondent biunivoquement aux
sous-catégories pleines de $C$ telles que la relation $F\in \ob R$
pour un objet de $C$ soit de nature locale; cette notion est donc
essentiellement indépendante de l'univers choisi $𝒰$. Enfin,
en vertu de ce qui a été dit dans \hyperref[IV.subsec8.2]{8.2}, une relation
de nature locale sur un topos (ou sur un site) est essentiellement
la même chose qu'un ouvert dudit topos (ou du site).

\setcounter{subsection}{3}
\subsection{Exemples d'ouverts}\etiquette[8.4]{IV.subsec8.4}

\subsubsection{}\etiquette[8.4.1]{IV.subsubsec8.4.1}
Soit $X$ un espace topologique, et interprétons $\Top(X)$
(\hyperref[IV.subsec2.1]{2.1}) comme la catégorie des espaces étales sur $X$.
Comme il est clair que les monomorphismes dans $\Top(X)$ sont les
applications injectives, il s'ensuit que les ouverts du topos $T(X)$
s'identifient aux ouverts de l'espace $X$. Plus généralement, les
ouverts d'un topos $\Top(X,G)$ (\hyperref[IV.subsec2.4]{2.4}) s'identifient
aux ouverts de $X$ \emph{invariants} par l'action de $G$.

\subsubsection{}\etiquette[8.4.2]{IV.subsubsec8.4.2}
Les ouverts d'un $𝒰$-topos $E$ forment un ensemble
$𝒰$-petit (\hyperrefext[I][I.prop7.4]{7.4}) ordonné par l'inclusion,
admettant des sup quelconques et des inf finis. Il admet l'objet
initial (ou « objet vide ») $\phi_E$ comme plus petit élément,
et l'objet final $e_E$ de $E$ comme plus grand élément. Ces deux
éléments ne sont identiques que si $E$ est un « topos vide »
(\hyperref[IV.subsec2.2]{2.2}).

\subsubsection{}\etiquette[8.4.3]{IV.subsubsec8.4.3}
Soit $G$ un Monoïde dans $E$, d'où un topos $B_G$
(\sisi{\hyperref[IV.subsec2.4]{2.4}}{\hyperref[IV.subsec2.6]{2.6}}), catégorie des objets de $E$ sur lesquels $G$
opère à gauche. L'objet final $e_G$ de $B_G$ est
l'objet\pageoriginale final $e_E$ de $E$ avec opération triviale de
$G$. On voit par suite que l'ensemble ordonné des ouverts de $B_G$
est canoniquement isomorphe à la catégorie des ouverts de $E$. En
particulier, si $E$ est le topos ponctuel donc $G$ est un monoïde
ordinaire, l'ensemble des ouverts de $B_G$ est exactement formé de
deux éléments, savoir $\phi_{B_G}$ et $e_G$.

\subsubsection{}\etiquette[8.4.4]{IV.subsubsec8.4.4}
Si $E$ est un topos de la forme $\widehat{C}$, où $C$ est une
$𝒰$-catégorie, alors (\hyperref[IV.subsec8.1]{8.1}) l'ensemble ordonné
des ouverts de $E$ est isomorphe à l'ensemble ordonné des cribles de
$C$.

\subsubsection{}\etiquette[8.4.5]{IV.subsubsec8.4.5}
Soit $X$ un espace topologique sobre non discret. Montrer que
$\TOP(X)$ (\hyperref[IV.subsec2.5]{2.5}) a des ouverts qui ne proviennent pas
d'ouverts de $X$, \ie d'ouverts de $\Top(X)$, par image inverse à
l'aide du morphisme canonique \eqhyperref[IV.eq4.10.1]{4.10.1} $\TOP(X)\rightarrow
\Top(X)$.

\subsection{Exemples de relations de nature locale}\etiquette[8.5]{IV.subsec8.5}

Pour tout objet $X$ du topos $E$, on désigne par $F\rightarrow F_X$
le \emph{foncteurs de localisation} $j^{\ast}_X:E\rightarrow
E_{/X}$, $Y\mapsto Y\times X$ (\hyperref[IV.subsec5.2]{5.2}).

\subsubsection{}\etiquette[8.5.1]{IV.subsubsec8.5.1}
Soit $u:F\rightarrow G$ un morphisme de $E$. La relation en
l'argument $X\in \ob E$,

$u_X:F_X\rightarrow G_X$ est un monomorphisme (\resp un
épimorphisme, \resp un isomorphisme)

\noindent est une relation de nature locale. En particulier, si $G$
est un Groupe de $E$, la relation « $G_X$ est le groupe unité sur
$X$ » est de nature locale; l'ouvert de $E$ qui lui correspond est
appelé le \emph{cosupport du Groupe $G$.}

\subsubsection{}\etiquette[8.5.2]{IV.subsubsec8.5.2}
Soient\pageoriginale $u$, $v:F\rightrightarrows G$ deux morphismes
de $E$. La relation en l'argument $X\in \ob E$
$$
u_X=v_X:F_X\longrightarrow G_X
$$
est une relation de nature locale. En particulier, si $G$ est un
Groupe de $E$, et $u$ une \emph{section} de $G$
(\hyperref[IV.subsubsec4.3.6]{4.3.6}), la relation en $X$
$$
\mbox{la section $u_X$ du Groupe $G_X$ de $E_{/X}$ est nulle}
$$
est de nature locale; l'ouvert de $E$ qui lui correspond est appelé
le \emph{cosupport de la section $u$ du Groupe $G$.}

\begin{enonce*}[remark]{Remarque 8.5.3}\etiquette[8.5.3]{IV.rem8.5.3}
Lorsque $E=\Top(X)$, le cosupport d'un faisceau en groupes $G$,
\resp d'une section d'un tel faisceau $G$, n'est autre que l'ouvert
de $X$ complémentaire du \emph{support de} $G$ \resp du
\emph{support de} $u$ dans la terminologie classique [TF].
\end{enonce*}

\setcounter{subsubsection}{3}
\subsubsection{}\etiquette[8.5.4]{IV.subsubsec8.5.4}
Soit $P(f)$ une relation en l'argument $f\in {\rm f\ell\,} C$, $C$
un site. Lorsque dans $C$ les produits fibrés sont représentables,
on dit que la relation $P(f)$ est \emph{stable par changement de
base (\resp stable par descente), \resp de nature locale) sur le
but} (ou \emph{sur la base}, ou « \emph{en bas} »), si pour tout
morphisme $f:X\rightarrow Y$ de $C$, la relation en l'argument
$Y'\in \ob C_{/Y}$:

« le morphisme $f':X'=\fprod{X}{Y'}{Y}\rightarrow Y'$ déduit de
$f$ par changement de base par le morphisme structural
$Y'\rightarrow Y$ satisfait $P(f')$ »

est stable par changement de base, \resp stable par descente, \resp
est de nature locale. Lorsque la topologie de $E$ est définie à
l'aide d'une prétopologie, on dit que la relation $P(f)$ est
\emph{de nature locale sur la source} (ou \emph{« en haut »}) si
pour tout morphisme $f:X\rightarrow Y$ de $C$ et toute famille
$(g_i:X_i\rightarrow X)_{i\in I}$\pageoriginale de $\Cov(X)$, la
relation $P(f)$ équivaut à la relation « $P(fg_i)$ pour tout $i\in
I$ ». Lorsque $\Cov(X)$ est formée de toutes les familles
couvrantes de $C$, cette condition signifie donc aussi que la
relation en l'objet $X'$ du site induit $C_{/X}$
\begin{center}
$P(fg)$, où $g:X'\rightarrow X$ est le morphisme structural
\end{center}
est de nature locale. Noter que si la relation $P(f)$ est de nature
locale en haut, elle est a fortiori de nature locale en bas.

\subsubsection{}\etiquette[8.5.5]{IV.subsubsec8.5.5}
Prenons par exemple $C=(\Sch)$, catégorie des schémas $\in
𝒰$, avec la topologie fidèlement plate quasi-compacte (\SGA
3 IV 6.3) ou une topologie moins fine. Alors chacune des propriétés
suivantes d'un morphisme est de nature locale en bas:

surjectif, radiciel, universellement ouvert, universellement finie,
propre, quasi-compact, quasi-compact et dominant, homéomorphisme
universel, séparé, quasi-séparé, localement de type fini, localement
de présentation finie, de type fini, de présentation finie, un
isomorphisme, un monomorphisme, une immersion ouverte, une immersion
fermée, une immersion quasi-compacte, affine, quasi-affine, fini,
quasi-fini, entier, plat, fidèlement plat, net, lisse, étale

\noindent (\EGA IV 2.6.4, 2.7.1 et \EGA IV 17.7.1). D'autre part,
munissons $(\Sch)$ de la prétopologie pour laquelle, pour tout
schéma $X\in 𝒰$, $\Cov(X)$ est formé des familles de
morphisme $(f_i:X_i\rightarrow X)_{i\in I}$ qui sont surjectives et
telles que les $f_i$ soient plats et localement de présentation
finie. Alors chacune des propriétés suivantes est \emph{de nature
locale en haut}:

localement de type fini, localement de présentation finie, de type
fini, plat, net, lisse, étale.

\noindent (\EGA IV 17.7.5, 17.7.7).

\begin{enonce*}[remark]{Exercice 8.6}\etiquette[8.6]{IV.exer8.6}
Soit\pageoriginale $f:E\rightarrow F$ un morphisme de
$𝒰$-topos. Considérons la relation en l'argument $X\in \ob
F$: le morphisme induit $E_{/f^{\ast}(X)}\rightarrow F_{/X}$ est une
équivalence de topos. Prouver que cette relation est de nature
locale.
\end{enonce*}

\begin{enonce*}[remark]{Exercice 8.7}\etiquette[8.7]{IV.exer8.7}
(\emph{Partitions d'un topos, somme de topos}). Soit $E$ un
$𝒰$-topos. Une famille $(e_i)_{i\in I}$ d'ouverts de $E$,
\ie de sous-objets de l'objet final $e$ de $E$, est appelée une
\emph{partition de} $e$ (ou encore, du topos $E$) si le morphisme
canonique $\amalg_{i\in I}e_i\rightarrow e$ est un isomorphisme, \ie
(\hyperrefext[II][II.cor4.6.2]{4.6.2})) si $e$ est le sup des $e_i$ et si $i\neq j$
implique $e_i\cap e_j=∅_E$.
\begin{enumerate}
\renewcommand{\theenumi}{\alph{enumi}}
\renewcommand{\labelenumi}{\theenumi)}
\item Montrer que pour que la famille $(e_i)_{i\in I}$ d'objets de $E$
soit une partition de $E$, il faut et il suffit que le foncteur
\begin{equation*}
E\longrightarrow \prod_i E_{/e_i}\tag{8.7.1}
\end{equation*}\etiquette[8.7.1]{IV.eq8.7.1}
défini par les foncteur de localisation $E\rightarrow E_{/e_i}$ soit
une équivalence de catégories.

\item Soit réciproquement $(E_i)_{i\in I}$ une famille de
$𝒰$-topos, avec $I\in 𝒰$. Prouver que $E=\Pi_{i\in
I} E_i$ est un $𝒰$-topos (qui sera appelé le \emph{topos
somme} de la famille des topos $(E_i)_{i\in I}$). Définir une
partition $(e_i)_{i\in I}$ de $E$ et des équivalences de catégories
$E_{/e_i}\approx E_i$, de telle façon que les foncteurs de
projection $E\rightarrow E_i$ s'identifient aux foncteurs de
localisation $E\rightarrow E_{/e_i}$.

\item Soit $E$ le topos somme de la famille des topos $(E_i)_{i\in
I}$. Montrer que pour tout $i\in I$ le foncteur projection
$E\rightarrow E_i$ est de la forme $u^{\ast}_i$, où
\begin{equation*}
u_i:E_i\longrightarrow E\tag{8.7.2}
\end{equation*}\etiquette[8.7.2]{IV.eq8.7.2}
est\pageoriginale un morphisme de topos. Prouver que pour tout
$𝒰$-topos $F$, le foncteur $v\mapsto (v\circ u_i)_{i\in I}$
\begin{equation*}
\Hhomtop(E,F)\longrightarrow \prod_i\Hhomtop(E_i,F)\tag{8.7.3}
\end{equation*}\etiquette[8.7.3]{IV.eq8.7.3}
est une équivalence de catégories.

\item Soit $(e_i)_{i\in I}$ une partition du topos $E$. Pour tout
morphisme de topos $g:F\rightarrow E$, la famille $(f_i)_{i\in I}$,
avec $f_i=g^{\ast}(e_i)$, est une partition de $F$, et les
morphismes induits $g_i:E_{/e_i}\rightarrow F_{/f_i}$ permettent de
reconstituer $g$ (à isomorphisme unique près), le foncteur
$g^{\ast}$ s'identifiant au produit cartésien des foncteurs
$g^{\ast}_i$ (compte tenu des équivalences du type (\hyperref[IV.exer8.7]{8.7}
1)). En conclure une description complète de la catégorie
$\Hhomtop(F,E)$, en termes des catégories de la forme
$\Hhomtop(F',E_i)$, où $F'$ est un topos de la forme $F_{/f}$ ($f$
sommande directe de l'objet final de $F$) et $E_i=E_{/e_i}$.

\item En particulier, si $F$ est connexe non vide (\cf
  \hyperref[IV.subsubsec4.3.5]{4.3.5}) \ie si
pour toute partition $(f_i)_{i\in I}$ de $F$ il existe un $i\in I$
et un seul tel que $f_i\neq \phi_F$, prouver que le foncteur
canonique
\begin{equation*}
\coprod_{i\in I}\Hhomtop(F,E_i)\longrightarrow
\Hhomtop(F,E)\tag{8.7.4}
\end{equation*}\etiquette[8.7.4]{IV.eq8.7.4}
déduit des morphismes de topos \eqhyperref[IV.eq8.7.2]{8.7.2} est une équivalence
de catégories. Plus particulièrement, la famille des morphismes de
topos \eqhyperref[IV.eq8.7.2]{8.7.2} induit une équivalence de catégories
$$
\coprod_{i\in I}\Ppoint(E_i)\xrightarrow{\approx}\Ppoint(E){\quoi}.
$$

\item Une\pageoriginale partie $I$ de l'ensemble des ouverts de $E$
est appelée une \emph{partition réduite} de $E$ si la famille
identique $(e_i)_{i\in I}$ indexée par $I$ est une partition, et si
les $e_i$ sont $\neq \phi_E$. Montrer que la relation de raffinement
(\hyperrefext[I][I.sec4.3.2]{4.3.2}) entre partitions réduites fait de l'ensemble $P$
des partitions réduites un ensemble ordonné $𝒰$-petit, tel
que la borne supérieure de deux éléments de $P$ existe. On trouve
ainsi un système projectif strict $(I_p)_{p\in P}$, qui sera
considéré comme un pro-ensemble et noté $\pi_{\circ}(E)$. On
l'appelle le \emph{pro-$\pi_{\circ}$ du topos} $E$.

\item Prouver que $\pi_{\circ}(E)$ pro-représente le foncteur
$$
T\longmapsto f_{\ast}f^{\ast}(T)=\Gamma(E,
T_E):(𝒰\text{-}\Ens)\longrightarrow (𝒰\text{-}\Ens)
$$
associé au morphisme canonique $f:E\rightarrow P$ de $E$ dans le
topos ponctuel $P$ (\hyperref[IV.subsec4.3]{4.3}). En particulier, pour que ce
foncteur soit représentable, il faut et il suffit que
$\pi_{\circ}(E)$ soit essentiellement constant, \ie isomorphe (en
tant que pro-ensemble) à un ensemble « ordinaire », qui sera
noté encore $\pi_{\circ}(E)$. Pour que $\pi_{\circ}(E)$ soit réduit
à un point, il faut et il suffit que $E$ soit « connexe non
vide ». Pour que $\pi_{\circ}(E)$ soit vide, il faut et il suffit que
$E$ soit le « topos vide »  (\hyperref[IV.subsec2.2]{2.2}).

\item Supposons $E$ quasi-compact, \ie que tout recouvrement de son
objet final $e$ admette un sous-recouvrement fini. Montrer qu'alors
$\pi_{\circ}(E)$ est un ensemble profini, et peut s'identifier par
suite (moyennant l'équivalence de catégories bien connue entre
pro-objets de la catégorie des ensembles finis, et espaces compacts
totalement discontinus) à un espace compact totalement discontinu,
qu'on notera également $\pi_{\circ}(E)$.

\item Soit $X$ un espace topologique, $\pi_{\circ}(X)$ l'ensemble des
composantes connexes de $X$, considéré comme un pro-ensemble
constant. Définir un morphisme\pageoriginale canonique de
pro-ensembles
\begin{equation*}
\pi_{\circ}(X)\longrightarrow
\pi_{\circ}(\Top(X)){\quoi},\tag{8.7.5}
\end{equation*}\etiquette[8.7.5]{IV.eq8.7.5}
ou ce qui revient au même, une application canonique
\begin{equation*}
\pi_{\circ}(X)\longrightarrow \varprojlim
\pi_{\circ}(\Top(X)){\quoi}.\tag{8.7.6}
\end{equation*}\etiquette[8.7.6]{IV.eq8.7.6}
Montrer que le premier morphisme est un isomorphisme si et seulement
si $X$ est homéomorphe à un espace somme d'espaces connexes (ce qui
est le cas en particulier si $X$ est localement connexe).

\item Supposons que $X$ soit un schéma quasi-compact. Prouver qu'alors
\eqhyperref[IV.eq8.7.6]{8.7.6} est bijectif.

\item Établir le « caractère fonctoriel »  en $X$ des morphismes
\eqhyperref[IV.eq8.7.5]{8.7.5} et \eqhyperref[IV.eq8.7.6]{8.7.6}, en précisant pour
commencer le sens de cette expression.


\item (Comparer \hyperref[IV.exer7.6]{7.6} i)). Montrer que les deux conditions
suivantes sur le
   $𝒰$-topos $E$ sont équivalentes: (i) Le morphisme canonique de
   $E$ dans le topos ponctuel est « essentiel »  (exercice \hyperref[IV.exer7.6]{7.6} a))
   \ie le foncteur $I\mapsto I_E$ de $(𝒰\text{-}\Ens)$ dans $E$ commute
   aux produits indexés par des ensembles $\in 𝒰$. (ii) Pour tout
   objet $X$ de $E$, le topos induit $E_{/X}$ a un
   $\pi_{\circ}(E_{/X})$ qui est un ensemble ordinaire, \ie $X$ est
   isomorphe à la somme d'une famille d'objets connexes de $E$. - On
   dit alors que $E$ est un \emph{topos localement connexe} (comparer \sisi{IV }{} \hyperref[IV.exer2.7.5]{2.7.5}).
\end{enumerate}
\end{enonce*}

\begin{enonce*}[remark]{Exercice 8.8}\etiquette[8.8]{IV.exer8.8}
(\emph{Morphismes dominants de topos}).
\begin{enumerate}
\renewcommand{\theenumi}{\alph{enumi}}
\renewcommand{\labelenumi}{\theenumi)}
\item Soit $f:E'\rightarrow E$ un morphisme de topos. Montrer
l'équivalence des conditions suivantes:

\begin{itemize}
\item[(i)] $f_{\ast}(\phi_{E'})=\phi_E$ (où $\phi_E$, $=\phi_{E'}$
désignent les objets initiaux de $E$, $E'$).

\item[(ii)] Pour\pageoriginale tout objet $X$ de $E$ qui est « non
vide »  \ie non isomorphe à $\phi_E$, $f^{\ast}(X)$ est un objet
non vide de $E'$.

\item[$\rm(ii')$] Comme (ii), avec $X$ un sous-objet de l'objet final
$e_E$ de $E$, \ie $X$ un ouvert du topos $E$.
\end{itemize}

On dit alors que le morphisme $f$ de topos est \emph{dominant}.

\item Soit $f:X'\rightarrow X$ une application continue d'espaces
topologiques, d'où un morphisme de topos
$\Top(f):\Top(X')\rightarrow \Top(X)$ (\hyperref[IV.subsec4.1]{4.1}). Montrer
que $\Top(f)$ est dominant si et seulement si $f$ est dominant, \ie
$f(X')$ est dense dans $X$.

\item Montrer que si $f:E'\rightarrow E$ est un morphisme « conservatif »  de topos (\hyperref[IV.subsubsec6.4.0]{6.4.0}), \ie tel que
$f^{\ast}$ soit fidèle, alors $f$ est dominant. Montrer que le
réciproque n'est pas nécessairement vraie, même pour un
morphisme de
localisation $E_{/U}\rightarrow E$, où $U$ est un ouvert du topos
$E$, ((Montrer que dans ce cas $f$ n'est conservatif que si $U$ est
l'objet final de $E$ (donc si $f$ est une équivalence de topos), et
utiliser l'exemple b).)

\item Soit $f:X'\rightarrow X$ une flèche dans un topos $E$. On dit
que $f$ est un \emph{morphisme dominant} dans le topos $E$ si le
morphisme correspondant des topos induits $E_{/X'}\rightarrow
E_{/X}$ \eqhyperref[IV.eq5.5.2]{5.5.2} est dominant. Montrer que cette
propriété ne
dépend que de l'image $f(X')$ de $X'$ dans $X$, en tant qu'élément
de l'ensemble ordonné $\ouv(X)$ des sous-objets de $X$. Montrer que
le morphisme $E_{/X'}\rightarrow E_{/X}$ est conservatif si et
seulement si $f$ est couvrant.
\end{enumerate}
\end{enonce*}

